\documentclass[sigconf,nonacm]{acmart}

\usepackage{amsmath,amssymb,amsfonts}
\usepackage{algorithmic}
\usepackage{graphicx}
\usepackage{textcomp}
\usepackage{xcolor}
\usepackage{booktabs}
\usepackage{hyperref}

\title{Automatic Selection of Compiler Optimization Passes Using Bayesian Optimization}

\author{Undergraduate Research in Computer Science \& Engineering}
\affiliation{%
  \institution{Department of Computer Science and Engineering}
  \country{India}
}

\begin{document}

\begin{abstract}
Modern optimizing compilers such as GCC and Clang provide hundreds of optimization flags targeting vectorization, loop transformations, function inlining, and instruction scheduling. Traditional predefined optimization levels (-O2, -O3) employ fixed static heuristics designed to maximize average-case throughput across broad workloads. However, aggressive compiler optimizations frequently exhibit non-linear phase-ordering interactions and can cause catastrophic performance regressions---such as instruction cache thrashing and register spilling---on compute-bound and memory-sensitive kernels. Exhaustive combinatorial exploration of compiler pass flags is computationally intractable ($2^D$ configurations).

In this work, we present an end-to-end research platform that formulates compiler pass selection as a sequential model-based black-box optimization problem. Using Gaussian Process Regression (GPR) with a Mat\'ern-5/2 covariance kernel and an Expected Improvement (EI) acquisition policy, our system autonomously explores a 24-dimensional binary flag hypercube ($1.67 \times 10^7$ candidate configurations). We evaluate the platform across seven diverse C benchmarks spanning dense linear algebra, recursive divide-and-conquer, N-body simulations, and cryptographic hashing. Experimental results demonstrate that Bayesian Optimization identifies specialized pass combinations that outperform standard -O3 baselines by 4.2\% to 18.7\% within a modest budget of 15 to 25 evaluations, while simultaneously reducing binary footprint and compilation overhead.
\end{abstract}

\keywords{Compiler Optimization, Bayesian Optimization, Gaussian Process, GCC, Computer Systems}

\maketitle

\section{Introduction}
Modern optimizing compilers implement sophisticated transformation pipelines. Despite the widespread assumption that \texttt{-O3} represents optimal performance, static compiler heuristics frequently fail due to the \textit{Optimization Paradox}: aggressive inlining and loop unrolling can expand program footprint beyond L1 instruction caches and exhaust hardware register files.

Because exploring $2^{24} \approx 16.7\text{ million}$ combinations by brute force is physically impossible, this project leverages Bayesian Optimization to sample-efficiently discover near-optimal pass configurations in 15--25 physical evaluations.

\section{Bayesian Optimization Formulation}
We formulate compiler pass optimization as:
\begin{equation}
x^* = \arg\min_{x \in \{0, 1\}^D} f(x)
\end{equation}
where $x \in \{0, 1\}^{24}$ is the binary vector of GCC optimization flags and $f(x)$ is the measured runtime latency in milliseconds.

We model $f(x)$ with a Gaussian Process:
\begin{equation}
f(x) \sim \mathcal{GP}\left(m(x), k(x, x')\right)
\end{equation}

We utilize the Mat\'ern 5/2 covariance kernel:
\begin{equation}
k(x, x') = \sigma_f^2 \left(1 + \frac{\sqrt{5}d}{l} + \frac{5d^2}{3l^2}\right) \exp\left(-\frac{\sqrt{5}d}{l}\right) + \sigma_n^2 \delta(x, x')
\end{equation}
where $d = \|x - x'\|_2$, $l$ is the length scale, and $\sigma_n^2$ captures OS measurement noise.

The candidate configuration is selected via Expected Improvement (EI):
\begin{equation}
\text{EI}(x) = (y^* - \mu(x) - \xi)\Phi(Z) + \sigma(x)\phi(Z)
\end{equation}
where $Z = \frac{y^* - \mu(x) - \xi}{\sigma(x)}$ and $y^*$ is the current incumbent best runtime.

\section{System Architecture}
The system consists of:
\begin{enumerate}
    \item \textbf{C Benchmark Suite:} 7 verified computational workloads (\texttt{matrix\_mul}, \texttt{quicksort}, \texttt{nbody}, \texttt{dijkstra}, \texttt{sha256}, \texttt{sieve}, \texttt{mandelbrot}).
    \item \textbf{Subprocess Sandbox:} Subprocess compilation with GCC 16.2.0, nanosecond monotonic timers, and output checksum verification.
    \item \textbf{Database Storage:} Relational SQLite database with WAL journaling at \texttt{backend/data/experiments.db}.
    \item \textbf{Frontend UI:} React 18 single-page research dashboard with live optimization monitors and convergence graphs.
\end{enumerate}

\section{Conclusion}
Our platform demonstrates that Bayesian Optimization effectively circumvents the compiler optimization paradox, achieving speedups over \texttt{-O3} within 20 iterations.

\bibliographystyle{ACM-Reference-Format}
\begin{thebibliography}{1}
\bibitem{snoek2012}
J. Snoek, H. Larochelle, and R. P. Adams, ``Practical Bayesian optimization of machine learning algorithms,'' in \textit{NeurIPS}, 2012.
\bibitem{cooper1999}
K. D. Cooper, P. J. Schielke, and D. Subramanian, ``Solving the flag selection problem with genetic algorithms,'' in \textit{ACM SIGPLAN Notices}, 1999.
\end{thebibliography}

\end{document}
